% ============================================================
% Section 171: 双势垒微扰下的能级差变化
% ============================================================
\section{量子力学：双势垒微扰下的能级差变化}
\label{sec:double-barrier-perturbation}

\begin{modelbox}[题目：双势垒微扰下的能级差变化]
一维无限深势阱在 $x=0$ 及 $x=L$ 处有两壁。两个宽为 $a$、高为 $V$ 的小微扰势垒分别位于 $L/4$ 和 $3L/4$ 处，$a\ll L/100$。利用微扰方法估计该微扰所产生的 $n=2$ 与 $n=4$ 能级间能量差的变化。
\end{modelbox}

\begin{tipbox}[考点快速分析]
\begin{itemize}
    \item \textbf{一级微扰}：$E_n^{(1)}=\langle\psi_n|H'|\psi_n\rangle$，即微扰势在未微扰态中的期望值
    \item \textbf{积分中值定理}：$a\ll L$ 时，被积函数在窄区间内近似常数
    \item \textbf{波节与波腹}：$n=2$ 的波腹恰在势垒位置；$n=4$ 的波节恰在势垒位置
    \item \textbf{能级差变化}：$\Delta(E_4-E_2)=E_4^{(1)}-E_2^{(1)}$
\end{itemize}
\end{tipbox}

\begin{derivationbox}[标准解题过程]
\textbf{1. 未微扰系统}

一维无限深势阱（宽度 $L$）的波函数和能级：
\[
\psi_n(x)=\sqrt{\frac{2}{L}}\sin\frac{n\pi x}{L},\quad E_n^{(0)}=\frac{\pi^2\hbar^2}{2\mu L^2}n^2.
\]

\textbf{2. 一级微扰能量修正}

微扰势 $H'$ 为位于 $x_1=L/4$ 和 $x_2=3L/4$ 处的两个方势垒（宽 $a$，高 $V$）。一级修正：
\[
E_n^{(1)}=\int_{L/4-a/2}^{L/4+a/2}\!V\,|\psi_n|^2dx+\int_{3L/4-a/2}^{3L/4+a/2}\!V\,|\psi_n|^2dx.
\]

因 $a\ll L$，用积分中值定理近似：
\[
E_n^{(1)}\approx\frac{2Va}{L}\Bigl[\sin^2\frac{n\pi}{4}+\sin^2\frac{3n\pi}{4}\Bigr].
\]

\textbf{3. $n=2$ 与 $n=4$ 的修正}

\textit{$n=2$}：$\sin^2(\pi/2)=1$，$\sin^2(3\pi/2)=1$，故
\begin{equation}
E_2^{(1)}=\frac{2Va}{L}(1+1)=\frac{4Va}{L}.
\end{equation}

\textit{$n=4$}：$\sin^2(\pi)=0$，$\sin^2(3\pi)=0$，故
\begin{equation}
E_4^{(1)}=\frac{2Va}{L}(0+0)=0.
\end{equation}

\textbf{4. 能级差的变化}

微扰前：$\Delta E^{(0)}=E_4^{(0)}-E_2^{(0)}=\dfrac{\pi^2\hbar^2}{2\mu L^2}(16-4)=\dfrac{6\pi^2\hbar^2}{\mu L^2}$。

微扰后：$\Delta E=E_4-E_2=\Delta E^{(0)}+(E_4^{(1)}-E_2^{(1)})$。

能级差的变化量：
\begin{equation}
\boxed{\delta(\Delta E)=E_4^{(1)}-E_2^{(1)}=-\frac{4Va}{L}.}
\end{equation}

即 $n=2$ 与 $n=4$ 的能级差减小（因 $E_2$ 上升而 $E_4$ 不变）。注意 $E_2^{(1)}>E_4^{(1)}$，因此 $E_4-E_2$ 变小。
\end{derivationbox}

\begin{formulabox}[答案汇总]
\begin{center}
\begin{tabular}{ll}
\toprule
\textbf{结论} & \textbf{结果} \\
\midrule
$E_2$ 一级修正 & $E_2^{(1)}=4Va/L$ \\
$E_4$ 一级修正 & $E_4^{(1)}=0$ \\
能级差变化 & $\delta(\Delta E)=-4Va/L$ \\
物理原因 & 势垒位于 $n=2$ 波腹、$n=4$ 波节处 \\
\bottomrule
\end{tabular}
\end{center}
\end{formulabox}

\begin{insightbox}[物理图像]
\begin{itemize}
    \item 本题是微扰论与波函数空间分布关系的经典示例：微扰对能级的影响正比于波函数在微扰区域的密度。波腹处影响最大，波节处影响为零。
    \item 该原理在半导体量子阱设计中直接应用：通过在有源区（波函数最大值处）引入势垒或势阱，可精确调控特定能级的能量。
    \item 若势垒位置偏离 $L/4$ 和 $3L/4$，则 $n=4$ 的一级修正也不再为零，能级差的变化将更为复杂。
\end{itemize}
\end{insightbox}
